\documentclass[11pt,a4paper]{article}
\usepackage[utf8]{inputenc}
\usepackage[T1]{fontenc}
\usepackage[spanish]{babel}
\usepackage{tikz}
\usetikzlibrary{shapes.geometric,shapes.symbols,arrows.meta,positioning,fit,backgrounds,calc,babel}
\usepackage{geometry}
\geometry{a4paper, margin=2cm}
\usepackage{parskip}
\usepackage{hyperref}
\usepackage{xcolor}

\title{\textbf{Documentación de la Arquitectura}\\ \Large Práctica 4 - Desarrollo y Orquestación de Microservicios}
\author{Imad Habib Jali (901421) \\ Jorge Bellido Lobera (903080)}
\date{\today}

\begin{document}

\shorthandoff{<>}
\maketitle

\section{Introducción}
El presente documento detalla la arquitectura del sistema de gestión de pedidos e inventario desarrollado para la Práctica 4. El sistema evoluciona desde una solución monolítica hacia un ecosistema de microservicios orquestados en Kubernetes. Se integran patrones de comunicación asíncrona (MOM), seguridad perimetral mediante API Key y el patrón de consistencia eventual \textbf{Saga} para garantizar la integridad de los datos entre MariaDB y PostgreSQL.

\section{1. Vista de Módulos (UML Package Diagram)}
Ilustra la estructura estática del código y las dependencias lógicas. Se destaca el uso de la librería compartida \textit{mom\_client.py} (evolución de la P3).

\begin{center}
\begin{tikzpicture}[
    package/.style={draw, rectangle, fill=blue!10, minimum width=3.2cm, minimum height=1.2cm, align=center, font=\sffamily\scriptsize, rounded corners=2pt},
    dependency/.style={-{Stealth}, dashed, thick}
]

\node[package] (gateway) {\guillemotleft module\guillemotright \\ \textbf{api-gateway}};
\node[package, below left=1.2cm and 0.5cm of gateway] (orders) {\guillemotleft module\guillemotright \\ \textbf{orders-service}};
\node[package, below right=1.2cm and 0.5cm of gateway] (inventory) {\guillemotleft module\guillemotright \\ \textbf{inventory-service}};
\node[package, below=3cm of gateway] (notif) {\guillemotleft module\guillemotright \\ \textbf{notifications-service}};
\node[package, left=1.5cm of orders, fill=orange!10] (mom) {\guillemotleft module\guillemotright \\ \textbf{mom\_client.py}};

\draw[dependency] (gateway) -- node[left, font=\tiny, sloped] {\guillemotleft proxy\guillemotright} (orders);
\draw[dependency] (gateway) -- node[right, font=\tiny, sloped] {\guillemotleft proxy\guillemotright} (inventory);
\draw[dependency] (orders) -- node[above, font=\tiny] {\guillemotleft rest\guillemotright} (inventory);
\draw[dependency] (orders) -- node[above, font=\tiny] {\guillemotleft import\guillemotright} (mom);
\draw[dependency] (notif) -| node[pos=0.2, right, font=\tiny] {\guillemotleft import\guillemotright} (mom);

\end{tikzpicture}
\end{center}

\newpage
\section{2. Vista de Componentes y Conectores (UML Component Diagram)}
Describe la interacción en tiempo de ejecución. El servicio de pedidos orquestas la transacción distribuida.

\begin{center}
\begin{tikzpicture}[
    component/.style={draw, rectangle, fill=green!10, text centered, text width=2.8cm, minimum height=1.3cm, font=\sffamily\scriptsize},
    port/.style={draw, rectangle, fill=white, minimum size=4pt, inner sep=0pt},
    connector/.style={-{Stealth}, thick},
    database/.style={draw, cylinder, shape border rotate=90, aspect=0.25, fill=yellow!20, minimum height=1.2cm, minimum width=1.6cm, align=center, text centered, font=\sffamily\tiny}
]

\node[component] (agw) {\guillemotleft component\guillemotright \\ \textbf{API Gateway}};
\node[component, below=1.5cm of agw] (ord) {\guillemotleft component\guillemotright \\ \textbf{Orders Service}};
\node[component, right=2.2cm of ord] (inv) {\guillemotleft component\guillemotright \\ \textbf{Inventory Service}};
\node[component, below=2cm of ord, fill=red!10] (brk) {\guillemotleft component\guillemotright \\ \textbf{Broker MOM}};
\node[component, right=2.2cm of brk] (not) {\guillemotleft component\guillemotright \\ \textbf{Notif Service}};

\node[database, left=1cm of ord] (db1) {\guillemotleft db\guillemotright \\ MariaDB};
\node[database, right=1cm of inv] (db2) {\guillemotleft db\guillemotright \\ PostgreSQL};

\draw[connector] (agw) -- node[right, font=\tiny] {HTTP/Auth} (ord);
\draw[connector] (ord) -- node[above, font=\tiny] {REST Async} (inv);
\draw[connector] (ord) -- node[left, font=\tiny] {TCP/JSON} (brk);
\draw[connector] (brk) -- node[above, font=\tiny] {Event Push} (not);
\draw[connector, dashed] (ord) -- (db1);
\draw[connector, dashed] (inv) -- (db2);

\end{tikzpicture}
\end{center}

\section{3. Vista de Despliegue (UML Deployment Diagram)}
Mapeo de los componentes sobre la infraestructura de Kubernetes (Pods y Volúmenes).

\begin{center}
\begin{tikzpicture}[
    nodeK8s/.style={draw, thick, rectangle, fill=gray!10, minimum width=10cm, minimum height=5cm, align=center, double distance=2pt},
    pod/.style={draw, rectangle, fill=white, text width=2.5cm, align=center, font=\sffamily\tiny, minimum height=0.8cm}
]

\node[nodeK8s] (cluster) at (0,0) {};
\node[font=\sffamily\small, anchor=north] at (cluster.north) {\guillemotleft device\guillemotright\ \textbf{Kubernetes Cluster (Minikube)}};

% Pods
\node[pod] (p1) at (-3, 1) {\guillemotleft pod\guillemotright \\ api-gateway};
\node[pod] (p2) at (0, 1) {\guillemotleft pod\guillemotright \\ orders-svc};
\node[pod] (p3) at (3, 1) {\guillemotleft pod\guillemotright \\ inventory-svc};
\node[pod] (p4) at (-3, -1) {\guillemotleft pod\guillemotright \\ broker-mom};
\node[pod] (p5) at (0, -1) {\guillemotleft pod\guillemotright \\ mariadb};
\node[pod] (p6) at (3, -1) {\guillemotleft pod\guillemotright \\ postgres};

% Volumes
\node[draw, cylinder, shape border rotate=90, aspect=0.2, fill=orange!20, minimum width=1cm, minimum height=0.6cm, font=\tiny, below=0.2cm of p4] (v1) {PVC};
\node[draw, cylinder, shape border rotate=90, aspect=0.2, fill=orange!20, minimum width=1cm, minimum height=0.6cm, font=\tiny, below=0.2cm of p5] (v2) {PVC};
\node[draw, cylinder, shape border rotate=90, aspect=0.2, fill=orange!20, minimum width=1cm, minimum height=0.6cm, font=\tiny, below=0.2cm of p6] (v3) {PVC};

\draw[thick, <->] (p1) -- (p2);
\draw[thick, <->] (p2) -- (p3);
\draw[thick, ->] (p2) -- (p4);

\end{tikzpicture}
\end{center}

\section{4. Vista de Instalación (UML File/Artifact Structure)}
Estructura física del entregable organizado por responsabilidades.

\begin{center}
\begin{tikzpicture}[
    folder/.style={draw, rectangle, fill=yellow!20, font=\sffamily\tiny, minimum width=3cm, rounded corners=1pt},
    file/.style={font=\sffamily\tiny, align=left}
]
\node[folder] (root) {Raiz: \textbf{/Practica4/}};
\node[file, below right=0.5cm and 0.3cm of root.west, anchor=west] (f1) {- \guillemotleft folder\guillemotright\ \textbf{api-gateway/}};
\node[file, below=0.2cm of f1.west, anchor=west] (f2) {- \guillemotleft folder\guillemotright\ \textbf{orders-service/}};
\node[file, below=0.2cm of f2.west, anchor=west] (f3) {- \guillemotleft folder\guillemotright\ \textbf{inventory-service/}};
\node[file, below=0.2cm of f3.west, anchor=west] (f4) {- \guillemotleft folder\guillemotright\ \textbf{notifications-service/}};
\node[file, below=0.2cm of f4.west, anchor=west] (f5) {- \guillemotleft folder\guillemotright\ \textbf{k8s/} (Manifiestos YAML)};
\node[file, below=0.2cm of f5.west, anchor=west] (f6) {- \guillemotleft artifact\guillemotright\ \textbf{arquitectura\_p4.pdf}};
\node[file, below=0.2cm of f6.west, anchor=west] (f7) {- \guillemotleft artifact\guillemotright\ \textbf{INSTRUCCIONES.md}};

\draw[thick, gray] (root.south west) ++(0.1,0) |- (f1.west);
\draw[thick, gray] (root.south west) ++(0.1,0) |- (f2.west);
\draw[thick, gray] (root.south west) ++(0.1,0) |- (f3.west);
\draw[thick, gray] (root.south west) ++(0.1,0) |- (f4.west);
\draw[thick, gray] (root.south west) ++(0.1,0) |- (f5.west);
\draw[thick, gray] (root.south west) ++(0.1,0) |- (f6.west);
\draw[thick, gray] (root.south west) ++(0.1,0) |- (f7.west);
\end{tikzpicture}
\end{center}

\section{5. Vista de Distribución (Arquitectura de Red)}
\begin{center}
\begin{tikzpicture}[
    box/.style={draw, rectangle, fill=gray!5, text width=2.5cm, align=center, font=\sffamily\tiny, minimum height=1cm},
    net/.style={thick, dashed, <->}
]
\node[box] (ext) {Cliente Externo \\ (Navegador/Curl)};
\node[box, right=1.5cm of ext, fill=blue!5] (ing) {K8s Ingress \\ (mis-aplicaciones.com)};
\node[box, right=1.5cm of ing] (mesh) {Service Mesh \\ (Internal Network)};

\draw[net] (ext) -- node[above, font=\tiny] {Internet} (ing);
\draw[net] (ing) -- node[above, font=\tiny] {Port Forward} (mesh);
\end{tikzpicture}
\end{center}

\end{document}
